\section{Introduction}\label{sec:intro}

Governments devote substantial resources to startup promotion, especially at the local level, where such policies are justified by expected gains in innovation, competition, and regional spillovers. Yet many successful startups do not remain independent competitors. Instead, they exit through acquisition by incumbents. This raises a basic question: when buyout exit is pervasive, does subsidizing startup entry create socially valuable competition, or does it instead expand the supply of firms built to sell out?

This paper studies entry subsidies in an environment in which prospective entrepreneurs choose whether to enter, how much effort to exert, and whether to pursue an independence-oriented or a buyout-oriented project. The former is stronger as a stand-alone business; the latter is more valuable under acquisition. After success, the entrant and the incumbent may attempt a merger, which is approved with probability \(q\). A government chooses an entry subsidy before entry takes place. In the baseline, \(q\) should be read as a reduced-form summary of exit realizability: merger policy is one interpretation, but not the only one. Throughout, entrepreneurial direction refers to the direction of project choice between stand-alone-oriented and buyout-oriented entrepreneurship.

The core mechanism is straightforward. A more permissive merger environment raises the private value of entrepreneurship by improving the entrepreneur's exit option, but it also changes the direction of entrepreneurial project choice. When acquisition becomes sufficiently attractive, entrepreneurs shift from independence-oriented projects toward buyout-oriented projects. Because a uniform entry subsidy does not affect project choice, it cannot correct that composition margin; it only scales entry at the privately preferred project margin. As a result, easier buyout exit need not strengthen the case for startup subsidies and can instead reduce the optimal subsidy.

A second key point is methodological. A merger is not always attempted after startup success; it is attempted only when expected merger gains cover filing and implementation costs. Welfare analysis must therefore respect the same behavioral margins that govern private equilibrium. Once this is imposed, merger policy affects subsidy design only when it changes actual behavior.

The paper makes three contributions. First, unlike the literature on acquisitions and merger policy \citep{cunningham2021killer,letina2024killer,dijk2024direction,dijk2025strategic}, which studies how acquisition prospects affect innovation strategy and competition, this paper studies subsidy design when buyout exit jointly shapes entry, effort, and project orientation. The analytical novelty is a behaviorally consistent welfare framework: merger attempt is endogenous, so the same attempt margin must enter both private continuation values and welfare accounting. Section \ref{sec:subsidy} therefore derives a primitive benchmark in which the main subsidy formulas and comparative statics are pinned down by product-market primitives together with the merger-attempt condition alone. Second, unlike the literature on entrepreneurial finance and public startup support \citep{lerner2020vcrole,lerner2022governmentincentives,bai2021public}, which typically asks whether public support increases startup formation or financing, this paper shows that stronger exit opportunities can raise private incentives while weakening the social case for subsidy by redirecting entrepreneurship toward buyout-oriented projects. Uniform entry subsidies therefore scale entry but do not correct project composition. Third, unlike place-based policy work that emphasizes strategic competition across jurisdictions or generic local spillovers, this paper derives a local--national \emph{incidence} wedge around that primitive benchmark. For a fixed project, the wedge reflects planner-specific incidence weights; what buyout exit adds is that merger permissiveness can change which project is privately selected, thereby changing the economically relevant local--national comparison. Appendices \ref{app:sufficientlocalover} and \ref{app:sufficientlocalunder} develop the economic content of that incidence comparison. The clawback exercise in Section \ref{sec:extensions} illustrates instrument separation rather than constituting the paper's main analytical contribution.

These results matter for a broad set of startup policies, including grants, tax credits, accelerators, multilevel startup support, and broader startup policy packages \citep{fazio2020rdtax,denes2023investortax,gonzalezuribe2018accelerators,manaresi2021startupact,lanahan2015multilevel}. The main implication is that policy should not ask only whether support increases startup formation. It must also ask what kinds of startups are being created and whether public support reinforces incentives to build independent competitors or acquisition targets. Section \ref{sec:extensions} returns to this point through an illustrative clawback exercise that shows how an additional instrument can target the composition margin without changing the paper's core welfare logic.

The remainder of the paper is organized as follows. Section \ref{sec:related} reviews the related literature. Section \ref{sec:model} presents the model. Section \ref{sec:private} characterizes private equilibrium project choice, effort, and entry. Section \ref{sec:subsidy} derives the primitive benchmark subsidy under buyout exit. Section \ref{sec:decentralized} adds local and national incidence extensions around that benchmark. Section \ref{sec:extensions} considers extensions and policy design. Section \ref{sec:conclusion} concludes.
